\sect{पौष-शुक्ल-पुत्रदा-एकादशी-माहात्म्यम्}
\label{sec:padma-pausha-shukla-putrada}

\ganapatyadiDhyanam
\padmaPuranam

\dnsub{अथ त्रिचत्वारिंशोऽध्यायः॥४३}

\uvacha{युधिष्ठिर उवाच}

\twolineshloka
{कथिता वै त्वया कृष्ण सफलैकादशी शुभा}
{कथयस्व प्रसादेन शुक्लपक्षस्य या भवेत्}% ॥१॥

\twolineshloka
{किं नाम को विधिस्तस्याः को देवस्तत्र पूज्यते}
{कस्मै तुष्टो हृषीकेशस्त्वमेव पुरुषोत्तमः}% ॥२॥

\uvacha{श्रीकृष्ण उवाच}

\twolineshloka
{शृणु राजन् प्रवक्ष्यामि शुक्ला पौषस्य या भवेत्}
{कथयामि महाराज लोकानां हितकाम्यया}% ॥३॥

\twolineshloka
{पूर्वेण विधिना राजन् कर्तव्यैषा प्रयत्नतः}
{पुत्रदा नाम नाम्ना सा सर्वपापहरा परा}% ॥४॥

\twolineshloka
{नारायणोऽधिदेवोऽस्या कामदः सिद्धिदायकः}
{नातः परतरा काचित् त्रैलोक्ये सचराचरे}% ॥५॥

\twolineshloka
{विद्यावन्तं यशस्वन्तं करोति च नरं हरिः}
{शृणु राजन् प्रवक्ष्यामि कथां पापहरां पराम्}% ॥६॥

\twolineshloka
{भद्रावत्यां पुरा ह्यासीत् पुर्यां राजा सुकेतुमान्}
{तस्य राज्ञस्तथा राज्ञी चम्पका नाम वर्तते}% ॥७॥

\twolineshloka
{पुत्रहीनेन राज्ञा च काले नीतो मनोरथैः}
{नैवाऽऽत्मजं नृपो लेभे वंशकर्तारमेव च}% ॥८॥

\twolineshloka
{तेनैव राज्ञा धर्मेण चिन्तितं बहुकालतः}
{किं करोमि क्व गच्छामि सुतप्राप्तिः कथं भवेत्}% ॥९॥

\twolineshloka
{न राष्ट्रे न पुरे सौख्यं लेभे राजा सुकेतुमान्}
{साध्व्या स्वकान्तया सार्धं प्रत्यहं दुःखितोऽभवत्}% ॥१०॥

\twolineshloka
{तावुभौ दम्पती नित्यं चिन्ताशोकपरायणौ}
{पितरोऽस्य जलं दत्तं कवोष्णमुपभुञ्जते}% ॥११॥

\twolineshloka
{राज्ञः पश्चान्न पश्यामो योऽस्माकं तर्पयिष्यति}
{इत्येवं संस्मरन्तोऽस्य दुःखिताः पितरोऽभवन्}% ॥१२॥

\twolineshloka
{न बान्धवा न मित्राणि नामात्याः सुहृदस्तथा}
{रोचयन्त्यस्य भूपस्य न गजाश्वाः पदातयः}% ॥१३॥

\twolineshloka
{नैराश्यं भूपतेस्तस्य नित्यं मनसि वर्तते}
{नरस्य पुत्रहीनस्य नास्ति वै जन्मनः फलम्}% ॥१४॥

\twolineshloka
{अपुत्रस्य गृहं शून्यं हृदयं दुःखितं सदा}
{पितृदेवमनुष्याणां नानृणत्वं सुतं विना}% ॥१५॥

\twolineshloka
{तस्मात् सर्वप्रयत्नेन सुतमुत्पादयेन्नरः}
{इह लोके यशस्तेषां परलोके शुभा गतिः}% ॥१६॥

\twolineshloka
{येषां तु पुण्यकर्तॄणां पुत्रजन्मगृहे भवेत्}
{आयुरारोग्यसम्पत्तिस्तेषां गेहे प्रवर्तते}% ॥१७॥

\twolineshloka
{पुत्रजन्म गृहे येषां लोकानां पुण्यकारिणाम्}
{पुण्यं विना न च प्राप्तिर्विष्णुभक्तिं विना नृप}% ॥१८॥

\twolineshloka
{पुत्राश्च सम्पदो वाऽपि निश्चयादिति मे मतिः}
{एवं विचिन्त्यमानोऽसौ न शर्म लभते नृपः}% ॥१९॥

\twolineshloka
{प्रत्यूषे चिन्तयद्राजा निशीथे चिन्तयत्ततः}
{स्वयमात्मविनाशं च चिन्तयामास केतुमान्}% ॥२०॥

\twolineshloka
{आत्मघाते दुर्गतिं च चिन्तयित्वा तदा नृपः}
{दृष्ट्वाऽऽत्मदेहं पतितमपुत्रत्वं तथैव च}% ॥२१॥

\twolineshloka
{पुनर्विचार्याऽऽत्मबुद्ध्या आत्मनो हितकारणम्}
{अश्वारूढस्ततो राजा जगाम गहनं वनम्}% ॥२२॥

\twolineshloka
{पुरोहितादयः सर्वे न जानन्ति गतं नृपम्}
{गम्भीरे विपिने राजा मृगपक्षिनिषेविते}% ॥२३॥

\twolineshloka
{विचचार तदा राजा वनवृक्षान् विलोकयन्}
{वटानश्वत्थबिल्वांश्च खर्जूरान् पनसांस्तथा}% ॥२४॥

\twolineshloka
{बकुलान् सप्तपर्णांश्च तिन्दुकांस्तिलकानपि}
{शालांस्तालांस्तमालांश्च ददर्श सरलान्नृपः}% ॥२५॥

\twolineshloka
{इङ्गुदी ककुभांश्चैव श्लेष्मातकनगांस्तथा}
{शल्लकान् करमर्दांश्च पाटलान्बदरानपि}% ॥२६॥

\twolineshloka
{अशोकांश्च पलाशांश्च शृगालाञ्शशकानपि}
{वनमार्जारमहिषान् शल्लकांश्चमरानपि}% ॥२७॥

\twolineshloka
{ददर्श भुजगान् राजा वल्मीकादर्धनिःसृतान्}
{तथा वनगजान् मत्तान् कलभैः सह सङ्गतान्}% ॥२८॥

\twolineshloka
{यूथपांश्च चतुर्दन्तान् करिणीयूथमध्यगान्}
{तान् दृष्ट्वा चिन्तयामास आत्मनः स गजान्नृपः}% ॥२९॥

\twolineshloka
{तेषां स विचरन्मध्ये राजा शोभामवाप ह}
{महदाश्चर्यसंयुक्तं ददृशे विपिनं नृपः}% ॥३०॥

\twolineshloka
{मार्गे शिवारुतान् शृण्वन्नुलूकविरुतं तथा}
{तांस्तानृक्षमृगान् पश्यन् बभ्राम वनमध्यतः}% ॥३१॥

\twolineshloka
{एवं ददर्श गहनं नृपो मध्यगते रवौ}
{क्षुत्तृड्भ्यां पीडितो राजा इतश्चेतश्च धावति}% ॥३२॥

\twolineshloka
{नृपतिश्चिन्तयामास संशुष्कगलकन्धरः}
{मया तु किं कृतं कर्म प्राप्तं दुःखं यदीदृशम्}% ॥३३॥

\twolineshloka
{मया वै तोषिता देवा यज्ञैः पूजाभिरेव च}
{तथैव ब्राह्मणा दानैस्तोषिता मिष्टभोजनैः}% ॥३४॥

\twolineshloka
{प्रजाश्चैव सदा कालं पुत्रवत् पालिता भृशम्}
{कस्माद्दुःखं मया प्राप्तमीदृशं दारुणं महत्}% ॥३५॥

\twolineshloka
{इति चिन्तापरो राजा जगामैवाग्रतो वनम्}
{सुकृतस्य प्रभावेन सरो दृष्टमनुत्तमम्}% ॥३६॥

\twolineshloka
{मीनसंस्पर्शमानं च पद्मैश्चापरशोभितम्}
{कारण्डैश्चक्रवाकैश्च राजहंसैश्च शोभितम्}% ॥३७॥

\twolineshloka
{मकरैर्बहुभिर्मत्स्यैरन्यैर्जलचरैर्युतम्}
{समीपे सरसस्तस्य मुनीनामाश्रमान् बहून्}% ॥३८॥

\twolineshloka
{ददर्श राजा लक्ष्मीवान् शकुनैः शुभशंसिभिः}
{दक्षिणं प्रास्फुरन्नेत्रमथ सव्येतरः करः}% ॥३९॥

\twolineshloka
{प्रास्फुरन्नृपतेस्तस्य कथयञ्शोभनं फलम्}
{तस्य तीरे मुनीन् दृष्ट्वा कुर्वाणन्नैगमं जपम्}% ॥४०॥

\twolineshloka
{हर्षेण महताऽऽविष्टो बभूव नृपनन्दनः}
{अवतीर्य हयात् तस्मान्मुनीनामग्रतः स्थितः}% ॥४१॥

\twolineshloka
{पृथक्पृथग्ववन्देऽसौ मुनींस्तान्शंसितव्रतान्}
{कृताञ्जलिपुटो भूत्वा दण्डवच्च पुनः पुनः}% ॥४२॥

\onelineshloka*
{प्रत्यूचुस्तेऽपि मुनयः प्रसन्ना नृपते वयम्}

\uvacha{राजोवाच}

\twolineshloka
{के भवन्तोऽत्र कथ्यन्तां का चाऽऽख्या भवतामपि}
{किमर्थं सङ्गता यूयं सत्यं वदत मेऽग्रतः}% ॥४३॥

\uvacha{मुनय ऊचुः}

\twolineshloka
{विश्वेदेवा वयं राजन् स्नानार्थमिह चागताः}
{माघो निकटमायात एतस्मात् पञ्चमेऽहनि}% ॥४४॥

\twolineshloka
{अद्य चैकादशी राजन् पुत्रदा नाम नामतः}
{पुत्रं ददात्यसौ विष्णुः पुत्रदा कारिणां नृणाम्}% ॥४५॥

\uvacha{राजोवाच}

\twolineshloka
{एष वै संशयो मह्यं पुत्रस्योत्पादने महान्}
{यदि तुष्टा भवन्तो वै पुत्रं मे दीयतां तदा}% ॥४६॥

\uvacha{मुनिरुवाच}

\twolineshloka
{अद्यैव दिवसे राजन् पुत्रदा नाम वर्तते}
{एकादशीति विख्यातं क्रियतां व्रतमुत्तमम्}% ॥४७॥

\twolineshloka
{अभिषेकात् ततोऽस्माकं केशवस्य प्रसादतः}
{अवश्यं तव राजेन्द्र पुत्रप्राप्तिर्भविष्यति}% ॥४८॥

\twolineshloka
{इत्येवं वचनात् तेषां कृतं राज्ञा व्रतोत्तमम्}
{मुनीनामुपदेशेन पुत्रदा या विधानतः}% ॥४९॥

\twolineshloka
{द्वादश्यां पारणं कृत्वा मुनीन्नत्वा पुनः पुनः}
{आजगाम गृहं राजा राज्ञी गर्भमथाऽऽदधौ}% ॥५०॥

\twolineshloka
{पुत्रो जातः सूतिकाले तेजस्वी पुण्यकर्मणा}
{पितरं तोषयामास प्रजापालो बभूव सः}% ॥५१॥

\twolineshloka
{तस्माद् राजन् प्रकर्तव्यं पुत्रदा व्रतमुत्तमम्}
{लोकानां तु हितार्थाय तवाग्रे कथितं मया}% ॥५२॥

\threelineshloka
{एकचित्तास्तु ये मर्त्याः कुर्वन्ति पुत्रदा व्रतम्}
{पुत्रान् प्राप्येह लोके तु मृतास्ते स्वर्गगामिनः}
{पठनाच्छ्रवणाद् राजन्नग्निष्टोमफलं लभेत्}% ॥५३॥

॥इति श्रीपाद्मे महापुराणे पञ्चपञ्चाशत्साहस्र्यां संहितायामुत्तरखण्डे उमापतिनारदसंवादान्तर्गतकृष्णयुधिष्ठिरसंवादे पौष-शुक्ल-पुत्रदा-एकादशी-माहात्म्यं नाम त्रिचत्वारिंशोऽध्यायः॥४३॥

\genMangalaShloka

\hyperref[sec:ekadashi_mahatmyam_padma_puranam]{\closesub}
\clearpage

